\section{Methodology}
Our question is which protection strategy best contains an outbreak on the transport network
that connects Europe, and whether the answer depends on the type of disease. Answering it
fairly requires three components built in sequence: a transport network derived from open
data (\ref{building the map}), a metapopulation epidemic engine that runs on it (\ref{running outbreak}), and a set
of disease types and protection strategies evaluated on top (\ref{what to spread} and \ref{how to defend}). A
single principle constrains every component: across any comparison, only the factor under
study varies, while all other inputs are held fixed. The study is therefore comparative
rather than predictive. We do not estimate how many people a real epidemic would infect; we
quantify differences between scenarios (one strategy versus another, one disease type versus
another) under identical conditions.

\subsection{Building the map} \label{building the map}
The substrate is a weighted graph $G=(V,E)$ in which nodes $v\in V$ are cities and edges
$(i,j)\in E$ are travel routes carrying weight $w_{ij}$ proportional to traffic. We construct
$G$ with NetworkX from four open sources: air routes from OpenFlights, city populations from
GeoNames, road and rail topology from OpenStreetMap and GRIP, and ferry routes from
OpenStreetMap. Region is a parameter of construction, which is what allows us to hold the
construction \textbf{method} fixed while varying only geography.

Because a single air layer misrepresents how an epidemic actually propagates, the same node
set carries up to three layers, air, land, and water, combined into one shared-node multilayer
network (Figure 2). Edge weights are layer-specific and serve as passenger-volume proxies, but
they come from different recipes according to what data exists. Air weights are real route
frequencies from OpenFlights, and water weights are real ferry-route frequencies counted per
city pair from an OpenStreetMap ferry dump. That dump is filtered to passenger-carrying
services (OSM tag `passenger` not equal to `no`) whose endpoints are at least 5 km apart: the
first filter removes freight-only crossings, since the model concerns human movement, and the
distance threshold discards short river and canal crossings while retaining genuine straits
and sea ferries that connect distinct cities. Land has no comparably uniform relatively easy to flow data across
regions, so its weights are instead modelled with a radiation kernel
$K_{ij} = \mathrm{flux}_{ij}/\mathrm{pop}_i$ (Simini et al., 2012) computed from city
populations. The constraint we impose for fairness is not that all layers share one recipe,
but that each layer is built by a single method applied identically to every region
(OpenFlights for air, the same global ferry dump for water, the radiation kernel for land), so
that differences between regions reflect topology rather than the underlying data source;
European commuting data (Eurostat) is reserved for validating the land kernel within Europe,
not for parameterising any region. Node populations $N_i$ are taken from GeoNames where
available, with a degree-based proxy $N_i = P_0 + d(i)\,P_\mathrm{route}$ as a fallback.


\subsection{Running an outbreak on the map} \label{running outbreak}

Each city runs a compartmental model internally, while individuals move between cities along
the edges of $G$ once per day. This is a metapopulation reaction–diffusion process (Colizza
et al., 2007): a local \textbf{reaction} updates the compartments within each city, and a \textbf{mobility}
step couples cities through the network. Within city $i$ the force of infection is
$\lambda_i = \beta\, I_i / N_i$, and a daily step applies the reaction to every city and then
the mobility coupling before advancing to the next day.

The mobility step is where the multilayer map earns its keep. Treating all travel as a single
diffusive process is a known source of error (Balcan \& Vespignani, 2011), so the engine
implements two mechanisms matched to the layers:

\begin{itemize}
    \item \textbf{Diffusive mobility (air, water).} A fraction of each compartment leaves city $i$ for
  neighbour $j$ each day, proportional to the layer travel rate $\tau_\ell$ and the normalised
  edge weight $w_{ij}$, and mixes into the destination population. This is the standard
  reaction–diffusion coupling and carries a global invasion threshold in mobility rate
  (Colizza \& Vespignani, 2008).

  \item \textbf{Recurrent mobility (land, commuting).} Commuters travel to a neighbour by day, mix
  there, and return home at night, so they are not relocated. We model this by coupling the
  force of infection rather than moving population: residents of $i$ experience
  $\pi_i = \sum_j C_{ij}\, I^{*}_j / N^{*}_j$, where $C$ is the per-capita commuting matrix
  (the land kernel with a stay-home diagonal) and $I^{*}_j = \sum_i C_{ij} I_i$,
  $N^{*}_j = \sum_i C_{ij} N_i$ are the infectious and total people present at $j$. When no
  land layer is present $C$ is the identity and $\pi_i = I_i/N_i$, recovering the plain
  diffusive model.
\end{itemize}

The two mechanisms are distinct in principle: diffusive travel relocates individuals between cities, whereas recurrent commuting returns them to their home city after each day. Recurrent coupling is therefore the more appropriate representation of daily commuting behaviour.

To assess the impact of this modelling choice, we compared both mechanisms on the European air+land network (SIR with $\beta=0.32$, $\gamma=0.12$; air travel rate 0.0002; land commuting rate 0.3; seed size 2500; 210 days; averaged over seeds 0--2). The resulting epidemic trajectories were very similar, with peak active infection of approximately 71 million under recurrent mobility and 73 million under diffusive mobility. This similarity arises because the air layer already synchronises much of the network. Larger differences appear when the air layer is removed or on controlled chain networks, where diffusion produces earlier and higher peaks, although the effect remains below approximately 20%.

We therefore adopt recurrent mobility on modelling grounds rather than because it substantially changes outcomes in the baseline scenario. As a validation check, both approaches produce peak prevalence slightly below the theoretical peak expected for a fully synchronised SIR epidemic at this reproduction number (approximately 26\% of the population), indicating that the engine behaves within expectations.

Additionally, an optional extension allows transmission to occur during travel by applying the reaction dynamics to travellers for the duration of a journey. This effect is expected to be most relevant for long ferry crossings.


\subsection{Choosing what to spread} \label{what to spread}
The engine runs whichever compartmental model it is given, which makes the choice of diseases
the substantive decision. Rather than commit to named pathogens, whose exact parameters are
endlessly contestable, we test five disease \textbf{types}, each defined by its dynamical structure
and parameterised from one representative exemplar using values drawn entirely from the
literature (Table 1). The five types span the qualitative behaviours relevant to the question:
an immunizing acute epidemic (SIR), an infection with a latent period (SEIR), a lethal
infection controllable by isolation (SEIQR with a fatal compartment), an endemic infection
conferring no immunity (SIS), and a recurrent infection with waning immunity (SEIRS). Each
type adds one mechanism relative to the previous, so a difference in outcome is attributable
to that mechanism rather than to several simultaneous changes.

**Table 1: Disease types, models, and parameters.**

INSERT GO TO APPENDIX X TO SEE SOURCES

\subsection{Choosing how to defend} \label{how to defend}
With the disease types fixed, the remaining variable is the intervention, and the network
structure of Section 3.1 defines the available choices, because each intervention is a
decision about which structural elements to remove. Under a fixed budget we compare two modes.
The first is node targeting (vaccination): a fixed number of cities are rendered immune so
that they no longer transmit. The targeting rule is the object of study, and we compare no
intervention (control), random selection, selection by degree (the highest-traffic cities),
selection by betweenness (the structural bridges), and a structural rule based on $k$-core and
motif membership rather than raw connectivity (Kitsak et al., 2010). We additionally vary
coverage between low and high adherence. The second mode is edge targeting (interdiction): we
close routes rather than immunise cities, in four scenarios: (A) no intervention; (B) close
all air routes while retaining land and water; (C) close all air routes within an air-only
model; and (D) close only the top-$k$ routes ranked by degree versus by betweenness.
Scenarios B and C exist specifically to contrast the conclusion a single-layer model reaches
with the behaviour of the full multilayer network.

\subsection{Implementation and reproducibility} \label{implementation and reproducibility}
The study was implemented in Python using NumPy and NetworkX. The framework comprises three modules: data retrieval, network generation, and epidemic evaluation. Simulations are fully determined by their configuration and random seed, ensuring reproducibility across runs.

Most experiments were conducted through an interactive web-based simulator developed as part of the project. The simulator allows users to configure outbreak scenarios and observe epidemic progression through live epidemic curves and geographic visualisations, making it suitable both for experimentation and demonstration. A batch execution mode was also retained for running larger parameter sweeps and repeated trials.

Both interfaces use the same epidemic engine, configuration system, and output pipeline, ensuring that results are consistent regardless of how simulations are launched. The overall system architecture is shown in Figure x.
